\documentclass[11pt]{article}
\usepackage[margin=0.8in]{geometry}
\usepackage{amsmath,amssymb}
\usepackage[hidelinks]{hyperref}
\setlength{\parindent}{0pt}
\setlength{\parskip}{5pt}
\newcommand{\relu}{\operatorname{ReLU}}
\begin{document}
\begin{center}
{\Large\bfseries Backpropagation with Negative ReLU Inputs}\par
{\large One iteration: manual gradients and PyTorch verification}
\end{center}

\textbf{Source and scope.} The supplied \texttt{[HW\_Code2]} notebook is titled
\emph{DNN Back Propagation with Negative Values at ReLU}.
We calculate gradients at its initial weights using one input--target pair
and one backward pass, before any weight update. Vectors are row vectors;
$(W_\ell)_{ij}$ connects input $i$ to output $j$. Both layers have no bias.

\section*{1. Initial values and forward pass}
\[
X=\begin{bmatrix}1&2\end{bmatrix},\quad
T=\begin{bmatrix}0&0\end{bmatrix},\quad
W_1=\begin{bmatrix}2&1\\0&-1\end{bmatrix},\quad
W_2=\begin{bmatrix}-1&2\\1&2\end{bmatrix}.
\]
The hidden-layer preactivation and activation are
\begin{align*}
S_2=XW_1
 &=\begin{bmatrix}1(2)+2(0)&1(1)+2(-1)\end{bmatrix}
 =\begin{bmatrix}2&-1\end{bmatrix},\\
Y_2=\relu(S_2)&=\begin{bmatrix}2&0\end{bmatrix}.
\end{align*}
The output-layer preactivation and activation are
\begin{align*}
S_3=Y_2W_2
 &=\begin{bmatrix}2(-1)+0(1)&2(2)+0(2)\end{bmatrix}
 =\begin{bmatrix}-2&4\end{bmatrix},\\
Y_3=\relu(S_3)&=\begin{bmatrix}0&4\end{bmatrix}.
\end{align*}
For $s\ne0$, $\relu'(s)=1$ when $s>0$ and $0$ when $s<0$.
There are no zero preactivations here. The derivative masks are
\[
\relu'(S_2)=\begin{bmatrix}1&0\end{bmatrix},\qquad
\relu'(S_3)=\begin{bmatrix}0&1\end{bmatrix}.
\]

\section*{2. Loss and output error signal}
The default \texttt{nn.MSELoss()} averages over both output elements
\cite{mse}, so
\[
L=\frac{(y_{3,1}-t_1)^2+(y_{3,2}-t_2)^2}{2}
 =\frac{0^2+4^2}{2}=8.
\]
Consequently,
\[
\frac{\partial L}{\partial Y_3}
 =\frac{2}{2}(Y_3-T)=\begin{bmatrix}0&4\end{bmatrix},\qquad
\delta_3=\frac{\partial L}{\partial S_3}
 =\begin{bmatrix}0&4\end{bmatrix}\odot\begin{bmatrix}0&1\end{bmatrix}
 =\begin{bmatrix}0&4\end{bmatrix}.
\]
Here $\odot$ denotes elementwise multiplication. The first output ReLU
blocks its gradient because its input is negative; in this example its
output residual is also zero because $y_{3,1}=t_1=0$.

\section*{3. Gradient of the second-layer weights}
Since $s_{3,j}=\sum_i y_{2,i}(W_2)_{ij}$, the chain rule gives
$\partial L/\partial(W_2)_{ij}=y_{2,i}\delta_{3,j}$. Thus
\[
\boxed{\nabla_{W_2}L=Y_2^{\mathsf T}\delta_3
 =\begin{bmatrix}2\cdot0&2\cdot4\\0\cdot0&0\cdot4\end{bmatrix}
 =\begin{bmatrix}0&8\\0&0\end{bmatrix}.}
\]
The second row is zero because the second hidden activation is zero.

\newpage
\section*{4. Hidden error signal and first-layer gradients}
Before applying the hidden ReLU derivative, each hidden activation
receives the sum of its contributions to both output errors:
\[
\frac{\partial L}{\partial y_{2,i}}
 =\sum_j\delta_{3,j}(W_2)_{ij},\qquad
\frac{\partial L}{\partial Y_2}=\delta_3W_2^{\mathsf T}.
\]
Substituting the values,
\[
\delta_3W_2^{\mathsf T}
 =\begin{bmatrix}0&4\end{bmatrix}
  \begin{bmatrix}-1&1\\2&2\end{bmatrix}
 =\begin{bmatrix}0(-1)+4(2)&0(1)+4(2)\end{bmatrix}
 =\begin{bmatrix}8&8\end{bmatrix}.
\]
The second hidden preactivation is $-1$, so its ReLU derivative is zero:
\[
\delta_2=\frac{\partial L}{\partial S_2}
 =(\delta_3W_2^{\mathsf T})\odot\relu'(S_2)
 =\begin{bmatrix}8&8\end{bmatrix}\odot\begin{bmatrix}1&0\end{bmatrix}
 =\begin{bmatrix}8&0\end{bmatrix}.
\]
Since $s_{2,j}=\sum_i x_i(W_1)_{ij}$,
\[
\frac{\partial L}{\partial(W_1)_{ij}}=x_i\delta_{2,j},\qquad
\boxed{\nabla_{W_1}L=X^{\mathsf T}\delta_2
 =\begin{bmatrix}1\cdot8&1\cdot0\\2\cdot8&2\cdot0\end{bmatrix}
 =\begin{bmatrix}8&0\\16&0\end{bmatrix}.}
\]
Both weights feeding the second hidden neuron have zero gradient:
the negative preactivation blocks the nonzero upstream derivative $8$.

\section*{5. Verification with the executed notebook}
The notebook's \texttt{torch.matmul} implementation prints
\[
\texttt{W1.grad}=\begin{bmatrix}8&0\\16&0\end{bmatrix},\qquad
\texttt{W2.grad}=\begin{bmatrix}0&8\\0&0\end{bmatrix},\qquad
Y_3=\begin{bmatrix}0&4\end{bmatrix}.
\]
These exactly match the manual calculations, and the recorded loss is $8$.

\textbf{Accounting for \texttt{nn.Linear}.} PyTorch computes
$xA^{\mathsf T}+b$ and stores $A$ with shape (output features, input
features) \cite{linear}. For an equivalent network, set
\[
A_1=W_1^{\mathsf T}=\begin{bmatrix}2&0\\1&-1\end{bmatrix},\qquad
A_2=W_2^{\mathsf T}=\begin{bmatrix}-1&1\\2&2\end{bmatrix}.
\]
Since $(A_\ell)_{ji}=(W_\ell)_{ij}$, its parameter gradients are
\[
\nabla_{A_\ell}L=(\nabla_{W_\ell}L)^{\mathsf T},\qquad
\nabla_{A_1}L=\begin{bmatrix}8&16\\0&0\end{bmatrix},\quad
\nabla_{A_2}L=\begin{bmatrix}0&0\\8&0\end{bmatrix}.
\]
The added verification cell checks the manual matrices against both
direct gradients and both transposed \texttt{nn.Linear} gradients.
All four checks pass with maximum absolute error $0$. It also checks
both preactivations, both activations, and the loss against the values above.

\begingroup\small
\begin{thebibliography}{2}
\bibitem{linear} PyTorch documentation, \texttt{torch.nn.Linear}.
\url{https://docs.pytorch.org/docs/2.14/generated/torch.nn.Linear.html}
\bibitem{mse} PyTorch documentation, \texttt{torch.nn.MSELoss}.
\url{https://docs.pytorch.org/docs/2.14/generated/torch.nn.MSELoss.html}
\end{thebibliography}
\endgroup
\end{document}
